\documentclass[11pt]{article}
\usepackage[margin=1in]{geometry}
\usepackage{amsmath,amssymb,amsthm,hyperref}
\newtheorem{theorem}{Theorem}
\newtheorem{lemma}[theorem]{Lemma}
\newtheorem{corollary}[theorem]{Corollary}
\newcommand{\Q}{\mathbb Q}
\newcommand{\Z}{\mathbb Z}
\newcommand{\Res}{\operatorname{Res}}
\begin{document}
\setlength{\emergencystretch}{1em}
\title{The Gamma-shape resultant of consecutive binomial moments}
\author{}\date{}\maketitle
\begin{abstract}
We determine a universal divisor, exact degree, and positive leading
coefficient of the Gamma-shape resultant for consecutive univariate
binomial moments. The leading term follows from a finite-moment
Gaussian limit. A further exact divisor is available when the lower
exponent is one. On the positive real shape axis the resultant is
nonnegative, and any zero is repeated, so the remaining problem reduces
to the positive zeros of its repeated part. The resulting positivity
and squarefreeness routes remain open. An exact fixed-polynomial
certificate shows that Hurwitz stability of the normalized shape
resultant is already false for gap five.
\end{abstract}
Fix integers $a\ge0$, $d\ge1$, $n\ge1$. Put
\[
 P_j(t;\beta)=\sum_{k=0}^j\binom jk(aj+\beta)_{dk}t^k,
 \qquad R_j(u;\beta)=u^jP_j(u^{-1};\beta),
\]
and define
\[
 {\cal R}_{n,a,d}(\beta)=\Res_t(P_n,P_{n+1})
                       =\Res_u(R_n,R_{n+1}).
\]
The equality is the reciprocal-resultant identity with sign
$(-1)^{n(n+1)}=1$. It holds generically in $\beta$ and thus as a
polynomial identity, including specializations where $P_j$ loses degree.
The reciprocals $R_j$ are monic of exact degree $j$ at every shape.

\begin{theorem}[Exact resultant normalization]\label{thm:normalization}
The resultant is a nonzero polynomial in $\Z[\beta]$, with
\[
 \deg_\beta{\cal R}_{n,a,d}=(2d-1)\frac{n(n+1)}2,\qquad
 [\beta^{(2d-1)n(n+1)/2}]{\cal R}_{n,a,d}
 =d^{n(n+1)}\prod_{j=1}^n j^j .
\]
The monic polynomial
\[
 D_{n,a,d}(\beta)
 =\prod_{j=1}^n
    \bigl(\beta+a(n+1)\bigr)_{\max(dj-a,0)}
\]
divides $ {\cal R}_{n,a,d}$ in $\Z[\beta]$. Therefore
\[
 Q_{n,a,d}(\beta)=\frac{{\cal R}_{n,a,d}(\beta)}
                       {D_{n,a,d}(\beta)}
 \in\Z[\beta]
\]
has positive leading coefficient and degree
\[
 (2d-1)\frac{n(n+1)}2-\sum_{j=1}^n\max(dj-a,0).
\]
\end{theorem}
\begin{proof}[Proof of the forced divisor]
Fix $h\ge0$ with $a+h<dn$ and specialize
$\beta_0=-a(n+1)-h$. The two Gamma parameters are
\[
 an+\beta_0=-a-h,\qquad a(n+1)+\beta_0=-h.
\]
For an integer $m\ge0$, $(-m)_k=0$ exactly when $k>m$.
Thus $R_n(u;\beta_0)$ is divisible by
$u^{g_h}$, where $g_h=n-\lfloor(a+h)/d\rfloor>0$.
The other reciprocal is divisible by
$u^{n+1-\lfloor h/d\rfloor}$, whose exponent is at least $g_h+1$.
The specialized gcd therefore has degree at least $g_h$.

The Sylvester matrix at $\beta_0$ has corank at least $g_h$.
Constant invertible row operations over $\Q$ make at least $g_h$
rows vanish at $\beta_0$. Each such row is divisible by
$\beta-\beta_0$ before taking the determinant. Hence
$(\beta-\beta_0)^{g_h}$ divides the resultant.
The product of these distinct factors is
\[
 \prod_{h=0}^{dn-a-1}
 (\beta+a(n+1)+h)^{n-\lfloor(a+h)/d\rfloor}
 =
 \prod_{j=1}^n(\beta+a(n+1))_{\max(dj-a,0)} .
\]
If $dn\le a$, both products are empty. Since this divisor is monic,
integer coefficients of the quotient follow by monic division.
\end{proof}

\begin{lemma}[Finite-moment Gaussian limit]\label{lem:Gaussian}
For $\gamma>0$, let $\mathcal M_\gamma$ be the polynomial functional
$\mathcal M_\gamma(X^k)=(\gamma)_k$.
For fixed $m\ge1$, set $\gamma=\beta+am$ and
\[
 Y_{m,\beta}=\frac{X^d-\beta^d}{d\beta^{d-1/2}}.
\]
As $\beta\to+\infty$, every fixed moment
$\mathcal M_\gamma(Y_{m,\beta}^k)$ tends to the standard Gaussian
moment: zero for odd $k$, and $(k-1)!!$ for even $k$.
Only finitely many polynomial moments are used; no convergence of an
exponential generating function for $X^d$ is asserted.
\end{lemma}
\begin{proof}
For $Z_\gamma=(X-\gamma)/\sqrt\gamma$, the formal identity
\[
 \sum_{k\ge0}\mathcal M_\gamma(Z_\gamma^k)\frac{s^k}{k!}
 =e^{-s\sqrt\gamma}(1-s/\sqrt\gamma)^{-\gamma}
\]
has logarithm
\[
 \sum_{j\ge2}\frac{s^j}{j\,\gamma^{j/2-1}}.
\]
Each fixed coefficient therefore tends to that of $e^{s^2/2}$.
All fixed moments of $Z_\gamma$ converge and are bounded.
Since $\gamma/\beta\to1$ and
$(\gamma-\beta)/\sqrt\beta\to0$, this also holds for
$W=(X-\beta)/\sqrt\beta$.
Finally,
\[
 Y_{m,\beta}
 =W+\sum_{j=2}^d\frac1d\binom dj\beta^{(1-j)/2}W^j.
\]
A fixed power is a finite polynomial in $W$; every term beyond
$W^k$ has a coefficient tending to zero and a bounded moment.
\end{proof}

\begin{proof}[Proof of the degree and leading coefficient]
Let $\sigma=d\beta^{d-1/2}$ and define the monic polynomial
\[
 F_{m,\beta}(z)=\sigma^{-m}R_m(-\beta^d+\sigma z;\beta).
\]
Lemma~\ref{lem:Gaussian} gives coefficientwise convergence to
\[
 G_m(z)=\sum_{j=0}^{\lfloor m/2\rfloor}
        \frac{m!}{2^j j!(m-2j)!}\,z^{m-2j}.
\]
These obey $G_0=1$, $G_1=z$,
$G_{m+1}=zG_m+mG_{m-1}$. Consequently
\[
 J_m:=\Res(G_m,G_{m+1})
 =m^m\Res(G_m,G_{m-1})
 =m^mJ_{m-1},\qquad J_0=1.
\]
Exchanging arguments has sign $(-1)^{m(m-1)}=1$.
Hence $J_m=\prod_{j=1}^m j^j>0$.

The affine variable change and monic normalization give
\[
 {\cal R}_{n,a,d}(\beta)
 =\sigma^{n(n+1)}
       \Res_z(F_{n,\beta},F_{n+1,\beta}).
\]
Continuity of the resultant in its coefficients implies
\[
 \lim_{\beta\to+\infty}
 \frac{{\cal R}_{n,a,d}(\beta)}
      {\beta^{(2d-1)n(n+1)/2}}
 =d^{n(n+1)}J_n>0.
\]
The resultant is an integer polynomial in $\beta$. This finite
nonzero limit determines both its degree and leading coefficient.
\end{proof}

\begin{corollary}[The positive-shape target]\label{cor:positive-target}
For every $\beta>0$, $D_{n,a,d}(\beta)>0$, and consecutive
coprimality is equivalent to $Q_{n,a,d}(\beta)\ne0$.
For fixed $a,d,n$, nonvanishing of $Q$ on all positive $\beta$ is
equivalent to positivity there, since its leading coefficient is positive.
Coefficientwise nonnegativity of $Q$, or of its positive-content
primitive part, would suffice.
\end{corollary}
\begin{proof}
All factors of $D$ are positive on this interval. A real continuous
polynomial with positive leading coefficient and no positive zero
is positive on the whole positive half-line.
\end{proof}


\section{An additional divisor for lower exponent one}

\begin{theorem}[A forced complete collision at a negative shape]
\label{thm:lower-one-extra-factor}
For $a=1$ and all integers $d,n\ge1$, the polynomial
\[
 E_{n,d}(\beta)=\bigl(\beta+(d+1)n+d\bigr)^n
\]
divides $ {\cal R}_{n,1,d}(\beta)$.
It is relatively prime to $D_{n,1,d}$, so their product divides
the resultant in $\Z[\beta]$.
\end{theorem}
\begin{proof}
Put $\beta_* =-(d+1)n-d$ and $\gamma=-d(n+1)$.
The consecutive shapes are $\gamma$ and $\gamma+1$.
The identity
\[
 (\gamma+1)_{dk}=\frac{\gamma+dk}{\gamma}(\gamma)_{dk}
              =\frac{n+1-k}{n+1}(\gamma)_{dk}
\]
and
$\binom{n+1}{k}(n+1-k)/(n+1)=\binom nk$ show that
\[
 P_{n+1}(t;\beta_*)=P_n(t;\beta_*).
\]
The identity includes the vanishing top coefficient when $k=n+1$.
Thus $R_{n+1}(u;\beta_*)=uR_n(u;\beta_*)$, and the specialized
Sylvester corank is $n$. The argument used for the natural divisor
gives the asserted multiplicity.

Every zero of the natural divisor has the form
$-(n+1)-h$ with $0\le h\le dn-2$; if $dn=1$ that divisor is one.
Its largest offset is therefore $(d+1)n-1$, smaller than the offset
$(d+1)n+d$ of $\beta_*$. The two divisors are relatively prime.
\end{proof}

\section{Positivity on the real axis and repeated factors}

\begin{theorem}[Nonnegative resultant and repeated positive zeros]
\label{thm:nonnegative}
For all the stated integers $a,d,n$ and every real $\beta>0$,
\[
 {\cal R}_{n,a,d}(\beta)\ge0.
\]
If it vanishes at $\beta_0>0$, the two reciprocal polynomials have
a common factor of degree at least two, and
\[
 {\cal R}_{n,a,d}(\beta_0)
 ={\cal R}_{n,a,d}'(\beta_0)=0.
\]
The same repeated-zero conclusion holds after division by any polynomial
divisor having no positive real zero.
\end{theorem}
\begin{proof}
For $\gamma>0$, the polynomial functional has the integral representation
\[
 \mathcal M_\gamma(h)
 =\frac{1}{\Gamma(\gamma)}
   \int_0^\infty h(x)x^{\gamma-1}e^{-x}\,dx .
\]
The monomial moments are $(\gamma)_k$, by integration by parts.
For every real $u$, a reciprocal of even degree is strictly positive,
since it is $\mathcal M_\gamma((u+X^d)^{2j})$.
For one of odd degree, its derivative is
$(2j+1)\mathcal M_\gamma((u+X^d)^{2j})>0$.
It therefore has exactly one real root, which is simple.

Write the pair as a monic odd polynomial $O$ and a monic even
polynomial $F$. Exchanging their resultant has sign one.
At the single real root of $O$, $F$ is positive.
Every conjugate nonreal pair of roots of $O$ contributes
$|F(\xi)|^2\ge0$ to the product formula for $\Res(O,F)$, counting
multiplicities. This proves nonnegativity.

A common root cannot be real, and its conjugate is a second common
root. Thus the specialized Sylvester corank is at least two.
The same constant-row-operation argument as before shows that the
shape resultant has a zero of multiplicity at least two at $\beta_0$.
Dividing by a polynomial nonzero at $\beta_0$ preserves that multiplicity.
\end{proof}

\begin{corollary}[A repeated-part reduction]\label{cor:repeated-part}
Let $D_-$ be any real polynomial divisor of $ {\cal R}_{n,a,d}$ with
no positive real zero, and put $C={\cal R}_{n,a,d}/D_-$.
Consecutive coprimality at every positive real shape is equivalent
to $\gcd(C,C')$ having no positive real zero.
In particular, squarefreeness of $C$ suffices.
One may take $D_-=D_{n,a,d}$, or $D_-=D_{n,1,d}E_{n,d}$ when $a=1$.
\end{corollary}
\begin{proof}
A positive zero of $C$ is repeated by Theorem~\ref{thm:nonnegative},
so is a zero of $\gcd(C,C')$. The reverse implication is immediate.
\end{proof}

This reduction does not prove squarefreeness of the remaining cofactor.
Removing its squarefull part by definition would discard exactly the
possible positive-shape obstruction, so that operation alone is not a
proof of coprimality.

\section{A failure of shape-parameter Hurwitz stability}

\begin{theorem}[An exact beta-variable obstruction]\label{thm:beta-hurwitz}
Let $C(\beta)$ be the positive-leading primitive part of
\[
 \frac{{\cal R}_{7,0,5}(\beta)}
      {\prod_{j=1}^7(\beta)_{5j}} .
\]
This polynomial has degree $112$ and strictly positive integer
coefficients, but its leading Hurwitz minor of order $84$ is negative.
Consequently it has a zero in the open right half-plane.
This is a statement about the shape variable $\beta$.
It gives neither a positive real zero nor a counterexample to
consecutive coprimality.
\end{theorem}
\begin{proof}
The exact certificate records the complete $113$ descending coefficients
and the positive integer content removed from the quotient.
Its equality to the displayed quotient is verified as a degree-$252$
polynomial identity: the full resultants agree at the $253$ distinct
integer shapes $1,\ldots,253$, using literal monic reciprocal
polynomials and exact integer resultants.
The degree bound and divisor degree are given by
Theorem~\ref{thm:normalization}. Positivity and the degree are checked
coefficient by coefficient.

For descending coefficients $c_0,\ldots,c_{112}$, with coefficients
outside this range set to zero, form
\[
 H_{84}=(c_{2j-i})_{i,j=1}^{84}.
\]
An exact integer determinant calculation gives $\det H_{84}<0$.
The decimal determinant has $7555$ digits; its full signed-value hash
is recorded and checked with the coefficient certificate.
A separate rational Routh elimination gives its first negative pivot
at row $84$ and agrees with the determinant sign.

The Routh--Hurwitz criterion in
Holtz--Khrushchev--Kushel~\cite[Theorem~1, p.~1]{HKK} requires all
these leading minors to be positive for strict stability.
In fact the negative minor excludes even the closed left half-plane:
if every zero of $C$ had nonpositive real part, then
$C(z+\epsilon)$ would be strictly stable for every $\epsilon>0$.
Its order-$84$ minor would be positive, and continuity as
$\epsilon\downarrow0$ would imply $\det H_{84}\ge0$, a contradiction.
\end{proof}

The coefficientwise positive result here is a complete assertion about
one explicit polynomial. It is not evidence promoted to an unrestricted
coefficientwise positivity theorem. The obstruction also differs from
a resultant viewed as a polynomial in the lower exponent at fixed shape.

\paragraph{Exact replay and scope.}
The checker verifies the retained identities on declared bounded inputs,
the formal Gaussian coefficient limits on finite moment controls, and
the complete fixed-polynomial Hurwitz certificate. Universal claims
above are justified by their written proofs, rather than by a finite
parameter search. Coefficientwise positivity and suitable squarefreeness
of the residual resultants remain open, as does unrestricted ordinary
binomial SFC.

\begin{thebibliography}{9}
\bibitem{HKK}
O.~Holtz, S.~Khrushchev and O.~Kushel,
\emph{Generalized Hurwitz matrices, generalized Euclidean algorithm,
and forbidden sectors of the complex plane},
\href{https://arxiv.org/abs/1506.07379v1}{arXiv:1506.07379v1},
June 2015. Theorem~1 is the classical Routh--Hurwitz criterion used here.
\end{thebibliography}
\end{document}
